% =============================================================
% GIAO DIỆN THỬ NGHIỆM - SAO CHÉP THIẾT LẬP TỪ TÀI LIỆU TOÁN
% Chỉ điều chỉnh phần phụ thuộc môn học và tương thích nội dung Hóa.
% =============================================================
\usepackage[vietnamese]{babel}
\usepackage{fontspec}
\usepackage{unicode-math}
\IfFontExistsTF{TeX Gyre Termes}{\setmainfont{TeX Gyre Termes}}{\setmainfont{Latin Modern Roman}}
\IfFontExistsTF{TeX Gyre Pagella Math}{\setmathfont{TeX Gyre Pagella Math}}{\setmathfont{Latin Modern Math}}

\usepackage{amsmath,mathtools}
\allowdisplaybreaks
\usepackage{enumerate}
\usepackage{tasks}
\settasks{label=\alph*),label-width=18pt,label-offset=4pt,item-indent=24pt,column-sep=1em,after-item-skip=1pt}
\usepackage{tikz}
\usetikzlibrary{math,through,calc,intersections,angles,quotes,shapes,shapes.geometric,arrows,patterns,snakes,matrix,chains,arrows.meta,decorations.shapes,decorations.fractals,decorations.markings,shadows,positioning,decorations.text,decorations.pathmorphing,shadings,fadings}
\usepackage{pgfplots}
\usepackage{pgfornament}
\usepgfplotslibrary{fillbetween}
\pgfplotsset{compat=1.18}
\usepackage[hidelinks,unicode]{hyperref}
\usepackage{currfile}
\usepackage[outline]{contour}
\usepackage{fontawesome}
\usepackage{longtable,multirow,makecell,array,tabularx}
\usepackage{diagbox}
\renewcommand{\tabcolsep}{3mm}
\newcolumntype{C}[1]{>{\centering\arraybackslash}p{#1}}
\newcolumntype{L}[1]{>{\raggedright\arraybackslash}p{#1}}
\usepackage{marginnote}
\def\noteleft{5}
\pgfmathsetmacro{\so}{\noteleft - 0.3}
\usepackage[top=\tren cm,bottom=\duoi cm,left=\trai cm,right=\phai cm,marginparwidth=\so cm,marginparsep=-\so cm]{geometry}
\renewcommand*{\marginfont}{\small}
\usepackage{fullwidth,varwidth}
\newenvironment{mynote}{\begin{fullwidth}[innermargin=\noteleft cm,width=\linewidth-\noteleft cm]}{\end{fullwidth}}
\usepackage{changepage}
\strictpagecheck
\usepackage{lastpage}
\usepackage{fancyhdr}
\usepackage[most]{tcolorbox}
\usepackage[explicit]{titlesec}
\usepackage{titletoc}

% Nội dung header/footer: giữ dữ liệu của tài liệu Hóa, dùng hình thức của bản Toán.
\providecommand{\ChanTrang}{TRUNG TÂM DẠY HỌC NGOÀI GIỜ NHÂN TRÍ SƠN}

% HEADER AND FOOTER STYLING
%--------------------------
\newcommand{\myfancynum}{% trên
\begin{tikzpicture}[remember picture,overlay]
	\fill[fill=cyan] ([yshift=-\tren cm+\topset cm]current page.north west) rectangle +(\paperwidth,0.5mm) coordinate (AA);
	\foreach \ii/\mauphu in {5/white,4.8/magenta,4/white,3.8/green,3/white,2.8/orange,2/white,1.8/cyan}{
		\fill[fill=\mauphu,draw=white] ([yshift=-2mm-0.5*\topset cm]AA)--++(-\ii,0)--++(0.5*\topset cm+2mm,0.5*\topset cm+ 2mm)--++(\ii cm,0)--cycle;}
	%----
	\checkoddpage\ifoddpage %nếu trang lẻ
	\def\kctrai{\trai} \def\kcphai{\phai} \def\ndtren{\leftmark}
	\else
	\def\kctrai{\phai} \def\kcphai{\trai} \def\ndtren{\rightmark}
	\fi
	%------
	\node[anchor=south west,text=magenta,inner sep=0pt] at ([yshift=-\tren cm+\topset cm+1mm,xshift=\kctrai cm]current page.north west) {\fontfamily{qag}\fontsize{9pt}{1pt}\selectfont \nouppercase{\ndtren}};
	\draw ([yshift=-\tren cm+\topset cm+1mm,xshift=-\kcphai cm]current page.north east) node[anchor=south east,text=magenta,inner sep=0pt] {\fontfamily{qag}\fontsize{8pt}{1pt}\selectfont \TenTruong};
	%--tròn trái
	\fill[cyan] (current page.north west) circle (\tren -\topset);
	\draw[white,line width=1pt] (current page.north west) circle (\tren cm -\topset cm-3pt);
	\draw (current page.north west)+(-45:0.5*\tren -0.5*\topset) node[white]{\fontfamily{qag}\fontsize{7pt}{1pt}\selectfont\bfseries\thepage};
\end{tikzpicture}%
}
\newcommand{\headdethi}{% trên đề thi
	\begin{tikzpicture}[remember picture,overlay]
		\fill[fill=cyan] ([yshift=-\tren cm+\topset cm]current page.north west) rectangle +(\paperwidth,0.5mm) coordinate (AA);
		\foreach \ii/\mauphu in {5/white,4.8/magenta,4/white,3.8/green,3/white,2.8/orange,2/white,1.8/cyan}{
			\fill[fill=\mauphu,draw=white] ([yshift=-2mm-0.5*\topset cm]AA)--++(-\ii,0)--++(0.5*\topset cm+2mm,0.5*\topset cm+ 2mm)--++(\ii cm,0)--cycle;}
		%----
		\checkoddpage\ifoddpage %nếu trang lẻ
		\def\kctrai{\trai} \def\kcphai{\phai} \def\ndtren{\TenKyThi}
		\else
		\def\kctrai{\phai} \def\kcphai{\trai} \def\ndtren{\LoaiDe~\thesection}
		\fi
		%------
		\node[anchor=south west,text=magenta,inner sep=0pt] at ([yshift=-\tren cm+\topset cm+1mm,xshift=\kctrai cm]current page.north west) {\fontfamily{qag}\fontsize{9pt}{1pt}\selectfont \nouppercase{\ndtren}};
		\draw ([yshift=-\tren cm+\topset cm+1mm,xshift=-\kcphai cm]current page.north east) node[anchor=south east,text=magenta,inner sep=0pt] {\fontfamily{qag}\fontsize{8pt}{1pt}\selectfont \TenTruong};
		%--tròn trái
		\fill[cyan] (current page.north west) circle (\tren -\topset);
		\draw[white,line width=1pt] (current page.north west) circle (\tren cm -\topset cm-3pt);
		\draw (current page.north west)+(-45:0.5*\tren -0.5*\topset) node[white]{\fontfamily{qag}\fontsize{7pt}{1pt}\selectfont\bfseries\thepage};
	\end{tikzpicture}%
}
%---------------------
\newcommand{\myfancyfoot}{% dưới
\begin{tikzpicture}[remember picture,overlay]
	\coordinate (AA) at ([yshift=\duoi cm-0.75*\topset cm-1mm]current page.south west);
	\checkoddpage\ifoddpage %nếu trang lẻ
		\def\kctrai{\trai} \def\kcphai{\phai}
	\else
		\def\kctrai{\phai} \def\kcphai{\trai}
	\fi
	%------
	\pgfmathsetmacro{\daisach}{\kctrai+1} 
	\fill[cyan] ([xshift=\daisach cm+10mm]AA) rectangle +(\paperwidth,0.5mm) coordinate (BB);
	%
	\fill[fill=cyan!20] ([yshift=1mm]AA)--++(\daisach cm,0) to[bend left=8mm] ++(3mm,-1.5mm)coordinate (Tr1)--++(0,-0.75)to[bend right=8mm]++(-3mm,1.5mm)--++(-\daisach,0)--cycle;
	\fill[fill=cyan] (Tr1) to[bend left=8mm]++(0.6,3mm)--++(0,-0.75)to[bend right=8mm]([yshift=-0.75cm]Tr1)--cycle;
	\draw[cyan,line width=0.5mm] ([yshift=1mm]AA)--++(\daisach cm,0) to[bend left=8mm] ++(3mm,-1.5mm);
	%
	\draw ([yshift=-0.25cm,xshift=-3pt]Tr1) node[anchor=east,text=magenta,inner sep=0pt] {\fontfamily{qag}\fontsize{8pt}{1pt}\selectfont{\bfseries \thepage}{\footnotesize /}{\bfseries\pageref{LastPage}}};
	%
	\node[anchor=north east,text=magenta,inner sep=0pt] (thinh) at ($([xshift=-\kcphai cm]AA)+(\paperwidth,-2pt)$) {\fontfamily{qag}\fontsize{10pt}{1pt}\selectfont\faFileText\, \ChanTrang};
	%--tròn phải
	\fill[cyan] (current page.south east) circle (\duoi -0.75*\topset);
	\draw[white,line width=1.5pt] (current page.south east) circle (\duoi cm -0.75*\topset cm-3pt);
\end{tikzpicture}%
}
%----------------------
\usepackage{changepage}
\strictpagecheck
\usepackage{lastpage}
\usepackage{fancyhdr,lastpage}
\pagestyle{fancy}
\fancyhf{}
\fancypagestyle{plain}{
	\fancyhead[LO,RE]{}
	\fancyfoot[LO,RE]{\myfancyfoot}
}
\fancyhead[LO,RE]{\myfancynum}
\fancyfoot[LO,RE]{\myfancyfoot}
\renewcommand{\footrulewidth}{0pt}
\renewcommand{\headrulewidth}{0pt}
%--------------4.2
\usepackage[most]{tcolorbox}
\colorlet{tcbcol@back}{tcbcolback}
\colorlet{tcbcol@frame}{tcbcolframe}
%-------------- Khung (Trong main này ko sd)
\newtcolorbox[auto counter]{khung4}[1]{enhanced, breakable,
	before skip=1mm,after skip=1mm,
	left=1mm,right=1mm,top=2mm,bottom=1mm,
	colframe=cyan!70!black,colback=yellow!15,colbacktitle=cyan!10,coltitle=red!70!black,
	,boxrule=0.4mm,
	attach boxed title to top center=
	{yshift=-0.1mm-\tcboxedtitleheight/2,yshifttext=2mm-\tcboxedtitleheight/2},
	varwidth boxed title*=-3cm,
	boxed title style={boxrule=0.3mm,
		frame code={ \path[tcb fill frame] ([xshift=-4mm]frame.west)
			-- (frame.north west) -- (frame.north east) -- ([xshift=4mm]frame.east)
			-- (frame.south east) -- (frame.south west) -- cycle; },
		interior code={ \path[tcb fill interior] ([xshift=-2mm]interior.west)
			-- (interior.north west) -- (interior.north east)
			-- ([xshift=2mm]interior.east) -- (interior.south east) -- (interior.south west)
			-- cycle;} 
	},
	fonttitle=\fontfamily{qag}\fontsize{10}{0}
	\bfseries,
	fontupper=\sffamily,
	title={#1}
}
%---------------Dạng toán
\newcounter{dang}\setcounter{dang}{0}
\renewcommand{\thedang}{\arabic{dang}}
%---Dạng 1
\newtcolorbox{dang}[1]{
	fonttitle=\fontfamily{qag}\bfseries,%fontupper=\itshape,
	colframe=teal,colback=yellow!10,
	sharp corners, breakable, halign title=center,%adjusted title=center, %canh giữa DẠNG
	before skip=2mm,after skip=2mm,
	left=2mm,right=2mm,top=0mm,bottom=0mm,
	boxrule=0.6pt,
	title={
		\color{white} \faFolderOpen\ Dạng~\thesection.\stepcounter{dang}\thedang.\ #1}
	\addcontentsline{toc}{subsection}{\it\sffamily \faFolderOpen\ Dạng~\thesection.\thedang: #1}
}
%---------------------------------------------------------------
% ĐỊNH NGHĨA SECTION. SUBSECTION, SUBSUBSECTION ... THEO Ý RIÊNG
%---------------------------------------------------------------
\usepackage[explicit]{titlesec} % để gọi #1
\usepackage{titletoc} % gói lệnh chứa cả titlesec và titletoc
%=====================================
\setcounter{secnumdepth}{4} %độ sâu
\renewcommand{\thechapter}{\arabic{chapter}}
\renewcommand{\thesection}{\thechapter.\arabic{section}}
\renewcommand{\thesubsection}{\Alph{subsection}}
\renewcommand{\thesubsubsection}{\arabic{subsubsection}}
%--------------Tròn
\newcommand{\tron}[1]{
	\begin{tikzpicture}[baseline=(A.base)]%
		\node[circle,draw=cyan,line width=0.5pt,fill=white,inner sep=2pt,outer sep=1pt] (A) {\color{white} #1};
		\node[circle,draw=none,fill=cyan,inner sep=1pt,outer sep=1pt] (A) {\color{white} #1};
	\end{tikzpicture}%
}
%============================Mục lục - Chapter*
%\makeatletter
\titleformat{name=\chapter,numberless}[display]
{\color{blue}\fontsize{23pt}{1pt}\fontfamily{qag}\selectfont\bfseries}
{}
{-1em}
{%
	\begin{tcolorbox}		
		[
		enhanced,
		frame empty,
		top=0.2cm, 
		bottom=0.2cm,
		colback=white,
		overlay={
			\fill[cyan] (frame.south west) rectangle ([yshift=-0.5cm]frame.north east);
			\fill[yellow!40,rounded corners=5pt,draw=black,line width=1pt] (frame.south west)--(frame.south east)--++(0,0.25)--++(-6,0)--([xshift=2cm]frame.north)--++(-4,0)--([xshift=6cm,yshift=0.25cm]frame.south west)--([yshift=0.25cm]frame.south west)--cycle;
			%---
			\fill[cyan,opacity=0.25] (frame.south west) rectangle ([yshift=-0.75cm]frame.south east);
			\fill[yellow!40,rounded corners=5pt,draw=black,line width=1pt,opacity=0.25] (frame.south west)--(frame.south east)--++(0,-0.25)--++(-6,0)--([xshift=2cm,yshift=-1cm]frame.south)--++(-4,0)--([xshift=6cm,yshift=-0.25cm]frame.south west)--([yshift=-0.25cm]frame.south west)--cycle; 
			%--- 
			\foreach \i/\mau/\bk in {-2/green/0.15,-1/violet/0.2,0/red/0.25,1/violet/0.2,2/green/0.15}{
				\fill[draw=violet,ball color=\mau]([yshift=-0.4cm,xshift=0.7*\i cm]frame.south) circle (\bk cm);
			}
			\node[xshift=0mm, yshift=0cm,scale=1,right=0cm] (tentron) at ([yshift=-0.4cm,xshift=0.7*2.1 cm]frame.south) {\pgfornament[height=0.5cm,color=red,opacity=1]{14}};
			\node[xshift=0mm, yshift=0cm,scale=1,left=0cm] (tentron) at ([yshift=-0.4cm,xshift=-0.7*2.1 cm]frame.south) {\pgfornament[height=0.5cm,color=red,opacity=1]{11}};
		}
		]
		\color{blue}\fontsize{23pt}{1pt}\fontfamily{qag}\selectfont\bfseries\centering\MakeUppercase{#1}
	\end{tcolorbox}%
	\begin{tikzpicture}[overlay,remember picture]
		\foreach \i in {0,1,...,22}{\draw[cyan,line width=2pt,opacity=0.2] 
			([xshift=0.3*\i cm]current page.north west)--++(-135:10cm)
			;}
	\end{tikzpicture}
}
[\vspace{1cm}]
%-------Đn section---------------------------
%--Biến đếm cho các tùy chọn BT
\newcounter{stt}
\setcounter{stt}{0}
%--
%\titlespacing{\section}{0cm}{0cm}{0cm}[0cm]
\titleformat{\section}
{	\setcounter{dang}{0}\setcounter{stt}{0}
	\fontsize{20pt}{10pt}\fontfamily{put}\selectfont\color{red}\centering\bfseries}
{
	M{\large ỤC }{\fontsize{25pt}{1pt}\selectfont\thesection}.
}
{0.2em}
{
	\MakeUppercase{#1}
}
[\vspace{0pt}]
%-------Đn subsection---------------------------
\newcommand{\dnsub}[1]{
	\begin{tikzpicture}[line join=round,line cap=round]%
		\node[inner sep=3pt,text=white,rounded corners=0pt,fill=blue] (muccon) at (0,0){\hspace*{2mm}#1\hspace*{0.75cm}}; %trái
		\fill[fill=orange] ([xshift=-3pt]muccon.north east) --([xshift=-3pt]muccon.east) 
		.. controls +(-90:0.2) and +(10:0.2) ..  ([xshift=-0.35cm]muccon.south east)
		--(muccon.south east) .. controls +(10:0.2) and +(-90:0.2) .. ([xshift=0.35cm]muccon.east)--([xshift=0.35cm]muccon.north east)--cycle ; %phải
	\end{tikzpicture}%
}
%\titlespacing{\subsection}{0cm}{0cm}{0cm}[0cm]
\titleformat{\subsection}
{\normalfont\fontsize{13pt}{13pt}\fontfamily{qag}\selectfont\bfseries}
{}
{-2.5mm}
{	
	\dnsub{\thesubsection\,--\,\MakeUppercase{#1}}
}
[]
%----------ĐN subsubsection-----------------------
%\titlespacing{\subsubsection}{0pt}{4mm}{1mm}[0cm]
\titleformat{\subsubsection}
{
	\fontsize{11.5pt}{9pt}\fontfamily{qag}\selectfont\bfseries\color{cyan}}
{
	\hspace*{-2mm}
	\begin{tikzpicture}[baseline=(char.base),line join=round,line cap=round]
		\node[signal, signal to=east, 
		text=white, fill=cyan] (char) at (0,0){\thesubsubsection.};
	\end{tikzpicture}%
}
{0.75em}
{#1}
[]
%================= Đn chương
%--------------------
\titlespacing{\chapter}{0cm}{1.15cm}{0.3cm}[0cm] %1: , 2: Trên, 3: dưới
\titleformat{\chapter}
{\fontsize{22pt}{6pt}\fontfamily{qag}\selectfont\bfseries\color{white}} %định dạng chung
{} %đánh số
{0cm} %khoảng cách số và nội dung
{
	\parbox{13.5cm}{\raggedright\MakeUppercase{#1}}
	\begin{tikzpicture}[remember picture,overlay,line join=round,line cap=round] 
		%=========footer
		\fill[top color=orange!20,bottom color=white] (current page.north west) rectangle ([yshift=-5.5cm]current page.north east);
		\begin{scope}
			\clip (current page.north west) rectangle ([yshift=-6cm]current 	page.north east);
			\draw[line width=1pt,white,step=1cm,rotate=45] (current page.north west) grid ([yshift=-6cm]current 	page.north east);
		\end{scope}
		%------------------------------------
		\checkoddpage\ifoddpage %nếu trang lẻ
			\def\sokc{\trai}
		\else
			\def\sokc{\phai}
		\fi
		%==============
		%-- tâm đặt--
		\coordinate (tam) at ([xshift=\sokc cm+2cm,yshift=-\tren cm-1.5cm]current page.north west);
		%-- nền dưới chương
		\fill[magenta,draw=black!50,line width=2pt,rounded corners=5mm] ([yshift=-2cm]tam) rectangle 	([xshift=-1cm,yshift=1cm]current page.north west);
		%--	tròn
		\node[circular drop shadow = {shadow scale = 1.03},minimum height=1.6cm, circle,fill=green!20, draw=white, line width=3pt,scale=1.75,text=red] 	(sochuong) at (tam){
			\fontfamily{put}\fontsize{30pt}{1pt}\selectfont %\thechapter %cố định kích thước hình tròn
		};
		%-- số
		\node[scale=1.75,text=red]  at (tam){
			\fontfamily{put}\fontsize{30pt}{1pt}\selectfont\thechapter
		};
		%-- Chương
		\draw[white,decoration={raise=0pt,text effects along path, text={C}{h}{u}{y}{ê}{n}{ }{đ}{ề},text align={center},text effects/.cd,text along path,every character/.style={font={\fontfamily{put}\fontsize{20pt}{1pt}\selectfont} }},decorate] ($(tam)+(200:1.7cm)$) arc (200:90:1.7);
		;
		%============================= nội dung chương
		\coordinate (tenbaichinh) at ([xshift=0.5cm,yshift=-1.25cm]sochuong.east); % điểm đặt chữ
		\def\n{0.25} %độ dày của bóng
		\def\gocnghieng{45} %nghiêng của bóng
		\def\x{1.2}\def\y{1.2} %co-giãn chữ
		\def\mauchu{red} %màu chữ
		\def\maubong{black}
		\foreach \i in {0,0.01,...,\n}{
			\pgfmathsetmacro\rr{\n-\i}
			\pgfmathsetmacro\domo{100*\i/\n}
			\node[anchor=south west,inner sep=2pt,text=\maubong!\domo!white,xscale=\x,yscale=\y] at ($(tenbaichinh)+(\gocnghieng:\rr)$){\parbox{13.5cm}{\raggedright\MakeUppercase{#1}}};
		}
		%------	
		\node[anchor=south west,inner sep=2pt,text=\mauchu,xscale=\x,yscale=\y] at (tenbaichinh) {\parbox{13.5cm}{\raggedright\MakeUppercase{#1}}};
	\end{tikzpicture}
}
[\vspace{20pt}]
%================== ĐN Part
\titlespacing{\part}{0cm}{0cm}{0cm}[0cm]
\titleformat{\part}[display]
{\bfseries}
{}
{0pt}
{
	\begin{tikzpicture}[remember picture, overlay]
		%----Tên phần
		\node[inner sep=0pt,scale=2,left] (tenphan) at ([xshift=-7cm]\textwidth,0) {\color{violet}\fontsize{15pt}{1pt}\fontfamily{qag}\selectfont\bfseries\MakeUppercase{\partname}};
		%-----Nội dung phần
		\node[inner sep=0pt,below left,scale=2.5] (partname) at ([xshift=0.2cm,yshift=-0.2cm]tenphan.south east) 
		{
			\parbox{5.5cm}{
				\fontsize{20pt}{20pt}\fontfamily{qag}\selectfont\bfseries\color{red}\raggedleft
				\MakeUppercase{#1}
			}
		};
		%----Nền phải
		\def\mauphan{cyan}
		\fill[magenta] ([xshift=0.25cm,yshift=0cm]partname.south east) rectangle ([xshift=7cm,yshift=0cm]tenphan.north east) node[midway,white,scale=2.5]{\fontsize{50pt}{1pt}\fontfamily{put}\selectfont\bfseries\thepart};
		%----Kẻ dưới
		\fill[\mauphan] ([xshift=6.8cm,yshift=-0.25cm]partname.south east) rectangle +(-1.2*\textwidth,-0.2);
		%----Trái và dưới
		\fill[cyan]
		(current page.north west)
		--([xshift=1cm]current page.north west)
		--([xshift=1cm,yshift=10cm]current page.south west)
		to[out=-90,in=180]([xshift=8cm,yshift=3cm]current page.south west)
		--([yshift=3cm]current page.south east)
		--(current page.south east)
		--(current page.south west)
		;
		%--
		\fill[cyan,opacity=0.5]
		([xshift=1.5cm,yshift=11.5cm]current page.south west)
		to[out=-90,in=180]([xshift=9cm,yshift=4cm]current page.south west)
		--([yshift=4cm]current page.south east)
		--([yshift=1.5cm]current page.south east)
		--([xshift=1.5cm,yshift=1.5cm]current page.south west)
		;
		%---trang trí
		\coordinate (BB) at ([yshift=-10 cm,xshift=3cm]current page.north west);
		\foreach \x in {1,...,50}{
			\pgfmathsetmacro{\r}{abs(random()*0.7*2)}
			\pgfmathsetmacro{\f}{30*\r)}
			\pgfmathsetmacro{\i}{abs(random()*15)}
			\pgfmathsetmacro{\j}{abs(random()*15)}
			\draw[fill=magenta!20,opacity=0.5, draw=blue] ([xshift=\i cm,yshift=-\j cm]BB) node[cyan]{\fontsize{\f}{1pt}\selectfont \x} circle (\r);
		}
	\end{tikzpicture}
}
[
\vspace{2cm}
\thispagestyle{empty}
]
%============================ 
\def\itemKN{\color{blue}\faCheckSquareO}
\def\itemCI{\color{violet}\faCheckCircleO}
%%======= Thiết lập labelitem, labelenumerate
%\renewcommand{\labelitemi}{\color{red}\faCheckSquareO}
\renewcommand{\labelitemi}{\color{violet}\faCheckCircleO}
\renewcommand{\labelitemii}{\color{violet}\bf ---}
\renewcommand{\labelitemiii}{\color{violet}\bf +}
\renewcommand{\labelenumi}{\alph{enumi})}
%\renewcommand{\labelenumii}{\color{blue}\bf\arabic{enumi}.\arabic{enumii}}
%============================
%============================
% Canh chỉnh mục lục chính
\setcounter{secnumdepth}{4} %Độ sâu đánh số
\setcounter{tocdepth}{2} %Độ sâu mục lục
\contentsmargin{0cm}
%~~~~~~~~~~~~~~~~~~~~~
\titlecontents{part}[0pc]
{\addvspace{10pt}%
	\color{red!70!black}\fontsize{18pt}{1pt}\fontfamily{qag}\selectfont\bfseries 
}%
{}
{}
{\hspace*{6pt}\titlerule\hspace*{6pt}\LARGE\bfseries \thecontentspage
	\begin{tikzpicture}[remember picture, overlay]
		\draw[fill=red!70!black, rounded corners=0pt] (2pt,0) rectangle (6,0.1pt);
	\end{tikzpicture}
}%
%~~~~~~~~~~~~~~~~~~~~~
\titlecontents{chapter}[7.25pc] % nd cách trái
{\addvspace{10pt}%
	\color{violet!70!black}\fontsize{14pt}{1pt}\fontfamily{put}\selectfont\bfseries
}%
{\contentslabel[\chaptertitlename\,\thecontentslabel.]{7.25pc}} %nhãn
{}
{\hspace*{6pt}\titlerule\hspace*{6pt}\bfseries\Large \thecontentspage
}%
%~~~~~~~~~~~~~~~~~~~~~
\titlecontents{section}[3.6pc]
{\addvspace{8pt}\bfseries\color{blue!70!black}}
{\fontsize{13pt}{18pt}\selectfont\sffamily\contentslabel[{\S\thecontentslabel}\,\,--]{3.5pc}}
{}
{\hfill
	\thecontentspage
}
[]
%~~~~~~~~~~~~~~~~~~~~~
\titlecontents{subsection}[5.8pc]
{\addvspace{1.0pt}\color{blue}}
{\fontsize{12pt}{15pt}\selectfont\sffamily\contentslabel[\tron{\thecontentslabel}]{2.4pc}}
{}
{{\tiny\dotfill}\thecontentspage}
[]
%~~~~~~~~~~~~~~~~~~~~~
%--------------------------
% ĐỊNH NGHĨA CÁC MÔI TRƯỜNG 
%--------------------------

% =============================================================
% BỔ SUNG RIÊNG CHO HÓA HỌC
% =============================================================
\usepackage{chemfig}
\setchemfig{atom sep=18pt}
\newcommand{\cfig}{\chemfig}
\usepackage{xymtexpdf}
\usepackage[version=4]{mhchem}
\usepackage{chemmacros}
\RenewChemPhase\gas{k}

\usepackage{siunitx}
\sisetup{retain-explicit-plus,output-decimal-marker={,},exponent-product={\cdot},range-phrase=\text{--},list-units=single,list-final-separator=\text{ và },per-mode=symbol}
\DeclareSIUnit\mola{\mole\per\liter}
\DeclareSIUnit\Mola{M}
\DeclareSIUnit\rates{\mole\per\liter\per\second}
\DeclareSIUnit\ratem{mol\,L^{-1}\,phút^{-1}}
\DeclareSIUnit\rateh{mol\,L^{-1}\,giờ^{-1}}
\DeclareSIUnit\rconstsec{\liter\per\mole\per\second}
\DeclareSIUnit\klr{\gram\per\cubic\centi\metre}
\DeclareSIUnit\klrl{\gram\per\milli\liter}
\DeclareSIUnit\cpm{\joule\per\kelvin\per\mole}
\DeclareSIUnit\cps{\joule\per\kelvin\per\gram}
\DeclareSIUnit\dc{\degreeCelsius}
\DeclareSIUnit\pm{\pico\metre}
\DeclareSIUnit\nm{\nano\metre}
\DeclareSIUnit\hou{giờ}
\DeclareSIUnit\minu{phút}
\DeclareSIUnit\day{ngày}
\DeclareSIUnit\yea{năm}
\DeclareSIUnit\amu{u}
\DeclareSIUnit\kjm{\kilo\joule\per\mole}
\DeclareSIUnit\kv{\kelvin}
\DeclareSIUnit\ml{\milli\liter}
\DeclareSIUnit\pc{\percent}
\DeclareSIUnit\atm{atm}
\DeclareSIUnit\barunit{bar}
\DeclareSIUnit\mmHg{mmHg}
\DeclareSIUnit\ev{\electronvolt}
\DeclareSIUnit\mps{\metre\per\second}
\DeclareSIUnit\kpa{\kilo\pascal}
\DeclareSIUnit\ton{tấn}
\DeclareSIUnit\hectare{ha}
\newcommand{\dd}{\text{dd}}
\newcommand{\kt}{\text{kết tủa}}
\newcommand{\du}{\text{dư}}
\newcommand{\HieuSuat}{H}
\newcommand{\Cpercent}{C\%}
\newcommand{\CM}{C_{\mathrm M}}

% Toggle
\newif\ifHienLoiGiaiVD
\newif\ifHienNguon
\newif\ifHienDapSo
\newif\ifHienLoiGiaiBT
\HienLoiGiaiVDtrue
\HienNguontrue
\HienDapSotrue
\HienLoiGiaiBTfalse
\newcommand{\BanHocSinh}{\HienLoiGiaiVDfalse\HienNguonfalse\HienDapSofalse\HienLoiGiaiBTfalse}
\newcommand{\BanGiaoVien}{\HienLoiGiaiVDtrue\HienNguontrue\HienDapSotrue\HienLoiGiaiBTtrue}

% Môi trường nội dung Hóa, dùng ngôn ngữ trình bày của bản Toán.
\newcounter{vidu}[section]
\renewcommand{\thevidu}{\arabic{vidu}}
\newtcolorbox[use counter=vidu]{vidu}[1][]{enhanced jigsaw,breakable,arc=0mm,colback=pink!3,boxrule=0.6pt,colframe=red!90!black,drop lifted shadow=orange,left=2mm,right=2mm,top=2mm,bottom=2mm,before skip=2mm,after skip=2mm,title={VÍ DỤ \thevidu},fonttitle=\fontfamily{qag}\selectfont\bfseries,coltitle=red!90!black,colbacktitle=pink!3,#1}
\newtcolorbox{loigiaibox}[1][]{enhanced,breakable,sharp corners,colback=white,colframe=cyan,boxrule=.35pt,left=2mm,right=2mm,top=1.4mm,bottom=1.4mm,before skip=1mm,after skip=2mm,title={Lời giải},fonttitle=\fontfamily{qag}\selectfont\bfseries\color{blue},coltitle=blue,colbacktitle=white,titlerule=0pt,#1}
\NewDocumentEnvironment{loigiai}{+b}{\ifHienLoiGiaiVD\begin{loigiaibox}#1\end{loigiaibox}\fi}{}
\newtcolorbox{ghinho}[1][]{enhanced jigsaw,breakable,pad at break*=1mm,colback=green!5,boxrule=0pt,frame hidden,left=2mm,right=0pt,bottom=1.5pt,top=1.5pt,before skip=2mm,after skip=2mm,borderline west={1mm}{0cm}{green!70!black},title={Ghi nhớ},fonttitle=\bfseries,coltitle=green!40!black,#1}
\newtcolorbox{casiobox}[1][]{enhanced jigsaw,breakable,pad at break*=1mm,colback=yellow!10,boxrule=0pt,frame hidden,left=2mm,right=0pt,bottom=1.5pt,top=1.5pt,before skip=2mm,after skip=2mm,borderline west={1mm}{0cm}{orange},title={Sử dụng máy tính CASIO},fonttitle=\bfseries,coltitle=orange!70!black,#1}
\newtcolorbox{unitbox}[1][]{enhanced jigsaw,breakable,pad at break*=1mm,colback=green!5,boxrule=0pt,frame hidden,left=2mm,right=0pt,bottom=1.5pt,top=1.5pt,before skip=2mm,after skip=2mm,borderline west={1mm}{0cm}{cyan!70!black},title={Đổi đơn vị},fonttitle=\bfseries,coltitle=cyan!60!black,#1}
\newcounter{tuluyen}[chapter]
\renewcommand{\thetuluyen}{\arabic{tuluyen}}
\newenvironment{tuluyen}{\refstepcounter{tuluyen}\par\medskip\noindent\textbf{\color{blue}Bài \thetuluyen.}\enspace\ignorespaces}{\par\smallskip}
\NewDocumentEnvironment{loigiaiBT}{+b}{\ifHienLoiGiaiBT\begin{loigiaibox}#1\end{loigiaibox}\fi}{}
\newcommand{\dapsohoa}[1]{\ifHienDapSo\hspace{\fill}\mbox{}\linebreak[0]\hspace*{\fill}\mbox{\fontsize{8pt}{1pt}\selectfont\color{blue!80!black}{\color{red}{\faKey}}~#1}\fi}
\let\dapso\dapsohoa
\newcommand{\source}[1]{\ifHienNguon\par\smallskip{\footnotesize\color{gray}\textit{Nguồn/kiểu bài: #1}}\fi}
\newcommand{\skill}[1]{\tcbox[on line,boxsep=.5pt,left=2pt,right=2pt,top=1pt,bottom=1pt,colback=cyan!8,colframe=cyan!45!black]{\scriptsize\bfseries #1}}
\newcommand{\key}[1]{\tcbox[on line,boxsep=.2pt,left=2pt,right=2pt,top=.6pt,bottom=.6pt,colback=black!4,colframe=black!55,arc=1pt]{\sffamily\scriptsize #1}}
\newcommand{\casioseq}[1]{\texttt{#1}}
\newcommand{\rxn}[1]{\[\ce{#1}\]}
\newcommand{\given}[1]{\textbf{Dữ kiện:} #1}
\newcommand{\ask}[1]{\textbf{Yêu cầu:} #1}
\newcommand{\result}[1]{\hfill\textbf{Kết quả: }#1}
\usepackage{etoc}
\newcommand{\chaptertoc}{\localtableofcontents\bigskip}
\newtcolorbox{thuthach}[1][]{enhanced,breakable,colback=white,colframe=violet,boxrule=.5pt,title={Bài toán thử thách},#1}
\setlength{\headheight}{15pt}
\setlength{\parindent}{0pt}
